\documentclass[11pt]{article}
\usepackage[T1]{fontenc}
\usepackage{lmodern,microtype}
\usepackage[margin=1in]{geometry}
\usepackage{amsmath,amssymb,amsthm,mathtools,booktabs,hyperref,xurl}
\newtheorem{theorem}{Theorem}[section]
\newtheorem{lemma}[theorem]{Lemma}
\newtheorem{proposition}[theorem]{Proposition}
\newtheorem{corollary}[theorem]{Corollary}
\theoremstyle{remark}\newtheorem{remark}[theorem]{Remark}
\newcommand{\Z}{\mathbb Z}\newcommand{\R}{\mathbb R}\newcommand{\N}{\mathbb N}
\newcommand{\PP}{\mathcal P}\newcommand{\tr}{\operatorname{tr}}
\newcommand{\vol}{\operatorname{covol}}
\newcommand{\ep}{\varepsilon}
\title{A sharp Gram-determinant law with a growing number of zero minors}
\author{Research note prepared for the ProveIt program}
\date{27 November 2026}
\begin{document}\maketitle
\begin{abstract}
For an integral matrix of size $2\times n$ and rank two, divide its Gram determinant by the square of the gcd of its maximal minors. We determine the asymptotic minimum of this normalized determinant when the number $z$ of vanishing maximal minors satisfies $n=o(z)$ and $z=o(n^2)$. The answer is $n^6/(729\Psi_*^2z^2)$, with the same sharp constant for an upper bound of $z$ and for exactly $z$ zero minors. Here $\Psi_*$ is the maximum of an explicit one-variable totient sum. The lower bound uses a squared truncated potential of primitive lattice directions and the quadratic multiplicity budget encoded by the zero minors. The construction has asymptotically parabolic multiplicities, and a four-square correction enforces the exact integer value of $z$. The proof is unconditional and does not assume a statement about zeros of an $L$-function.
\end{abstract}

\section{Definitions and the main result}
Let $A=[v_1\ \cdots\ v_n]\in\Z^{2\times n}$ have rank two. Its columns are indexed; identical, opposite, and zero columns are allowed. Put
\[
 D(A)=\gcd_{i<j}|\det(v_i,v_j)|,\qquad
 G(A)=\frac{\det(AA^{\mathsf T})}{D(A)^2},\qquad
 z(A)=\#\{\{i,j\}:i<j,\det(v_i,v_j)=0\}.
\]
Zeros may be included or omitted in the gcd, and $D(A)>0$ by the rank assumption. Let
\[
 E_{\le}(n,z)=\min_{z(A)\le z}G(A),\qquad
 E_{=}(n,z)=\min_{z(A)=z}G(A)
\]
when the specified class is nonempty. The first class is always nonempty for $n\ge2$. The exact class will be shown nonempty in our asymptotic range. All admissible $G(A)$ are positive integers by Cauchy--Binet after lattice normalization, so a nonempty class has a minimum.

For $T>0$ set $w_j=\varphi(j)/j$ and
\begin{align}
 \Phi(T)&=\frac4{3T}\sum_{1\le j\le T}w_j(1-j^2/T^2)^{3/2},\label{eq:phi}\\
 \Psi(T)&=\frac{16}{15T}\sum_{1\le j\le T}w_j(1-j^2/T^2)^{5/2},\qquad
 \Psi_* =\max_{T>0}\Psi(T).\label{eq:psi}
\end{align}
Terms at an integer endpoint vanish. Lemma~\ref{lem:maximum} proves that the maximum exists and is larger than $1/\pi$; the definition does not rely on a numerical maximization or a uniqueness conjecture.

\begin{theorem}[Sharp law, including exact degeneracy]\label{thm:main}
Let $n\to\infty$ and let $z=z(n)$ be nonnegative integers such that
\[
 \frac{z}{n}\longrightarrow\infty,\qquad \frac{z}{n^2}\longrightarrow0.
\]
Then the exact class is nonempty for all sufficiently large $n$, and
\[
 E_{\le}(n,z)\sim E_{=}(n,z)\sim
 \frac{n^6}{729\Psi_*^2z^2}.
\]
Equivalently, writing $Q=n+2z$,
\[
 E_{\le}(n,z)\sim E_{=}(n,z)\sim
 \frac{4n^6}{729\Psi_*^2Q^2}.
\]
\end{theorem}

This is a different asymptotic regime from the generic case $z=0$. The latter is treated in the source article by a first truncated potential, with $E_2(n)=\kappa(n-1)^4+O(n^2)$ and $\kappa=(64\Phi_*^2)^{-1}$. Allowing $z$ to grow faster than $n$ changes the optimal scale to $n^6/z^2$, and a squared potential determines its sharp coefficient. If $z$ is of order $n^{1+\alpha}$ with $0<\alpha<1$, the new exponent is $4-2\alpha$.

\section{Normalization and multiplicity budgets}
\begin{lemma}[Integral normalization]\label{lem:normalization}
The infima defining $E_{\le}$ and $E_=$ can equally be taken over unnormalized determinants $\det(CC^{\mathsf T})$ of integral rank-two matrices. A minimizing matrix for this latter problem generates $\Z^2$, and every nonzero column is primitive in $\Z^2$.
\end{lemma}
\begin{proof}
The lattice $L=\sum_i\Z v_i$ has index $D(A)$ in $\Z^2$: this is the rank-two instance of the determinantal-divisor formula, or follows by taking a basis matrix $B$ of $L$, writing $A=BC$, and applying Cauchy--Binet to an integer right inverse of $C$. Thus $C$ is integral, its columns generate $\Z^2$, $D(C)=1$, and
\[
 G(A)=\det(CC^{\mathsf T}).
\]
An invertible linear transformation preserves the zero pattern of the pair minors. Conversely $G(C)\le\det(CC^{\mathsf T})$ for every integral matrix $C$. This proves equality of the two minima. A minimizing matrix cannot generate a lattice of index greater than one, since taking coordinates in a basis of that lattice strictly decreases its positive determinant. If a nonzero column is $k u$ with integral primitive $u$ and $|k|>1$, replacing it by $u$ leaves the zero pattern unchanged and strictly decreases the sum of squared minors, since some other column is independent of it. The new columns still generate $\Z^2$, because their span contains the old span. This proves the last assertion.
\end{proof}

For a determinant-one lattice $\Lambda\subset\R^2$, let $\PP(\Lambda)$ contain one representative of each antipodal pair of nonzero primitive lattice vectors. A finite multiplicity function $m:\PP(\Lambda)\to\N\cup\{0\}$ describes every primitive configuration after signs are ignored.

\begin{lemma}[Isotropization]\label{lem:trace}
For every positive-definite real $2\times2$ matrix $M$,
\[
 \min_{g\in\operatorname{SL}_2(\R)}\tr(gMg^{\mathsf T})=2\sqrt{\det M}.
\]
Consequently a normalized minimizing configuration may be placed in a determinant-one lattice so that its sum of squared column lengths is exactly $2\sqrt{E}$.
\end{lemma}
\begin{proof}
The arithmetic--geometric mean inequality for the two eigenvalues gives the lower bound. Equality holds for $g=(\det M)^{1/4}M^{-1/2}$, which has determinant one.
\end{proof}

\begin{lemma}[Zero columns and repeated directions]\label{lem:budget}
If a primitive configuration has $k$ zero columns and nonzero-direction multiplicities $(m_v)$, then
\[
 n'=\sum_v m_v=n-k,\qquad
 \sum_v m_v^2=n+2z(A)-2kn+k^2.
\]
If $z(A)\le z=o(n^2)$, then $k\le2z/(n-1)=o(n)$ and $\sum m_v^2\le Q:=n+2z$.
\end{lemma}
\begin{proof}
The zero columns account for $k(n-k)+\binom{k}{2}$ zero pair minors. The other zero minors are exactly $\sum_v\binom{m_v}{2}$. Expanding gives the identity. Since $k(n-k)+\binom{k}{2}=k(2n-k-1)/2\ge k(n-1)/2$, the stated bound on $k$ follows. Finally $2kn-k^2\ge0$.
\end{proof}

\section{Uniform primitive-lattice potentials}
For $R>0$ define
\[
 P_s(\Lambda,R)=\sum_{v\in\PP(\Lambda)}(R^2-|v|^2)_+^s,\qquad s=1,2.
\]
After a rotation, a reduced basis of a determinant-one lattice has the form
\[
 (a,0),\quad(a\theta,a^{-1}),\qquad a>0,\quad |\theta|\le\tfrac12,
\]
where $a$ is a shortest nonzero lattice length and $a\le(4/3)^{1/4}$. Indeed, a shortest vector is primitive and extends to a basis; subtracting its nearest integral multiple from the other basis vector gives the shear bound, and shortestness gives $a^2\le a^2/4+a^{-2}$.
Choose the positive second coordinate as representative off the horizontal line. The primitive directions on row $j\ge1$ are
\[
 (a(u+j\theta),j/a),\qquad u\in\Z,\quad\gcd(u,j)=1,
\]
and there is one horizontal primitive direction $(a,0)$. Write $T=aR$.

\begin{lemma}[Uniform squared potential]\label{lem:uniform}
As $R\to\infty$,
\[
 P_2(\Lambda,R)=R^6\Psi(T)+O(R^5\log(2R)),
\]
with an absolute implied constant, uniformly over every determinant-one lattice, including $T\to0$ and $T\to\infty$.
For fixed $T>0$, uniformly for bounded shears $\theta$,
\begin{align*}
 P_1(\Lambda,R)&=R^4\Phi(T)+O_T(R^2),\\
 P_2(\Lambda,R)&=R^6\Psi(T)+O_T(R^4),\\
 \#\{v\in\PP(\Lambda):|v|<R\}&=O_T(R^2).
\end{align*}
The fixed-$T$ assertions remain true after excluding any fixed finite collection of primitive directions in the integer coordinates $(u,j)$, provided their lengths are $O_T(R)$.
\end{lemma}
\begin{proof}
On a row $j\le T$, put $\eta_j=(R^2-j^2/a^2)^{1/2}$. For $f(x)=(\eta_j^2-x^2)_+^2$,
\[
 \int_\R f(x)\,dx=\frac{16}{15}\eta_j^5,
 \qquad \operatorname{Var}(f)=2\eta_j^4\le2R^4.
\]
For any spacing $b>0$ and shift $c$, the sum of a compactly supported bounded-variation function on $b\Z+c$ differs from its integral divided by $b$ by at most its total variation. This follows by comparing each sampled value with the average on an adjacent interval and summing variations. Expand the coprimality condition by
\[
 \mathbf1_{\gcd(u,j)=1}=\sum_{d\mid\gcd(u,j)}\mu(d).
\]
The integral term on row $j$ is therefore $(w_j/a)(16/15)\eta_j^5$, while its error is $O(R^4\tau(j))$. Summing the integral terms gives precisely $R^6\Psi(aR)$. The row errors sum to $O(R^4T\log(2T))$, since $\sum_{j\le T}\tau(j)\le T(1+\log T)$ for $T\ge1$. The single horizontal direction contributes at most $R^4$. As $T\le(4/3)^{1/4}R$, the total error is $O(R^5\log(2R))$. If $T<1$ there are no nonhorizontal rows, and the same bound holds. This proves the uniform statement without a compactness assumption on lattice shapes.

For fixed $T$, there are boundedly many rows and residue classes. The same bounded-variation estimate gives an error $O_T(R^2)$ for $P_1$, since the variation of $(\eta^2-x^2)_+$ is $2\eta^2$, and $O_T(R^4)$ for $P_2$. The first integral is $4\eta^3/3$. Counting the primitive points in each row interval gives $O_T(R/a+1)=O_T(R^2)$ points. Excluding finitely many directions changes $P_s$ by at most $O_T(R^{2s})$ and the count by $O(1)$.
\end{proof}

\begin{lemma}[Existence of the scalar maximum]\label{lem:maximum}
The continuous function $\Psi$ tends to $1/\pi$ as $T\to\infty$, and satisfies $\Psi(32/3)>31839/100000>1/\pi$. Thus $\Psi_*$ is attained at a finite $T_0>1$. Every maximizing $T_0$ satisfies
\[
 \Phi(T_0)=\frac32\Psi_*.
\]
\end{lemma}
\begin{proof}
The identity $\varphi(j)/j=\sum_{d\mid j}\mu(d)/d$ gives
\[
 \sum_{j\le x}w_j=\frac6{\pi^2}x+O(\log(2x)).
\]
Partial summation against $(1-x^2/T^2)^{5/2}$ yields the limit because
\[
 \int_0^1(1-u^2)^{5/2}\,du=\frac{5\pi}{32}.
\]
For completeness, the strict inequality at $T=32/3$ is a small rational certificate. Put $s_j=1024-9j^2$ and $r_j=\lfloor10^{12}\sqrt{s_j}\rfloor$ for $1\le j\le10$. Then
\[
 \Psi(32/3)\ge\frac{16}{5\cdot32^6\cdot10^{12}}
 \sum_{j=1}^{10}\frac{\varphi(j)}{j}s_j^2r_j>\frac{31839}{100000}.
\]
The last comparison is between rational numbers and can be checked by integer square roots and cross multiplication. The classical bound $\pi>223/71$ gives $1/\pi<71/223<31839/100000$. Consequently the supremum exceeds the limiting value and is attained in a compact subinterval of $(1,\infty)$. The terms entering at integers vanish to order $5/2$, so $\Psi$ is differentiable there as well. Differentiating the finite sums gives
\[
 4\Phi(T)=6\Psi(T)+T\Psi'(T).
\]
At a maximizer $\Psi'(T_0)=0$.
\end{proof}

\section{The lower bound}
Let a minimizing matrix be normalized, made primitive, and isotropized using Lemmas~\ref{lem:normalization} and~\ref{lem:trace}; let $S$ be its sum of squared lengths. Write $n'=n-k$ and $Q_0=\sum m_v^2\le Q$. For every $R>0$,
\begin{align*}
 S&\ge n'R^2-\sum_v m_v(R^2-|v|^2)_+\\
  &\ge n'R^2-\sqrt{Q\,P_2(\Lambda,R)}.
\end{align*}
The second inequality is Cauchy--Schwarz. Take
\[
 R=\frac{2n'}{3\sqrt{Q\Psi_*}}.
\]
Because $z=o(n^2)$, Lemma~\ref{lem:budget} gives $n'\sim n$ and $R\to\infty$. Lemma~\ref{lem:uniform} gives, uniformly over the minimizing lattice,
\[
 P_2(\Lambda,R)\le R^6\bigl(\Psi_*+o(1)\bigr).
\]
It follows that
\[
 S\ge\left(1-o(1)\right)\frac{4(n')^3}{27Q\Psi_*}
   =\left(1-o(1)\right)\frac{4n^3}{27Q\Psi_*}.
\]
Since $S=2\sqrt{E_{\le}(n,z)}$, squaring proves
\[
 E_{\le}(n,z)\ge(1-o(1))\frac{4n^6}{729Q^2\Psi_*^2}.
\]
It automatically supplies the same lower bound for the exact class whenever that class is nonempty. In particular, concentrating multiplicity on a few directions, introducing zero columns, and allowing arbitrary anisotropy cannot evade the lower bound.

\section{A construction with a controlled unused budget}
The next lemma records enough precision to correct the number of zero minors exactly. Removing a fixed set of directions is included because it is needed for the correction.

\begin{lemma}[Rounded parabolic construction]\label{lem:construction}
Fix a maximizing $T_0$ from Lemma~\ref{lem:maximum}, and a fixed finite set $F$ of primitive integer directions. Suppose $N/Q\to0$, $Q/N^2\to0$, and $Q\equiv N\pmod2$. Write $Z=(Q-N)/2$. If $\eta\to0$ satisfies
\[
 \frac{Q/N^2+N/Q}{\eta}\longrightarrow0,
\]
then, for all sufficiently large $N$, there is a rank-two integral matrix with $N$ nonzero primitive columns, using no directions from $F$, such that its zero-minor count $Z'$ satisfies
\[
 0\le Z-Z'=O(\eta Q),
\]
and in a determinant-one lattice its squared-length sum is
\[
 S\le(1+o(1))\frac{4N^3}{27Q\Psi_*}.
\]
The implicit constant in the budget estimate may depend on $T_0,F$ but not on $N,Q,\eta$.
\end{lemma}
\begin{proof}
Choose
\[
 R^2=\frac4{9\Psi_*}\frac{N^2}{Q}(1+10\eta),\qquad a=T_0/R,
\]
and take the rectangular lattice with basis $(a,0),(0,a^{-1})$. For large $R$, this is in the reduced form used above. Let $f_v=(R^2-|v|^2)_+$ for primitive directions not in $F$, and put
\[
 P_1=\sum_v f_v,\quad P_2=\sum_v f_v^2,\quad
 c=N/P_1,\quad \widetilde m_v=c f_v.
\]
Then $\sum_v\widetilde m_v=N$. From Lemma~\ref{lem:uniform} and $\Phi(T_0)=3\Psi_*/2$,
\begin{align*}
 \sum_v\widetilde m_v^2
 &=N^2\frac{P_2}{P_1^2}
 =\frac4{9\Psi_*}\frac{N^2}{R^2}\bigl(1+O(R^{-2})\bigr)\\
 &=\frac{Q}{1+10\eta}\bigl(1+O(Q/N^2)\bigr).
\end{align*}
In particular the real squared-multiplicity sum lies between $Q-O(\eta Q)$ and $Q-5\eta Q$ once $N$ is large.

Round each $\widetilde m_v$ down, and then add one to sufficiently many of the finitely many entries with positive fractional part to make the sum exactly $N$. This is possible because $N-\sum_v\lfloor\widetilde m_v\rfloor$ is the integer sum of the fractional parts. Every change has absolute value at most one. There are $d=O(R^2)$ entries, and therefore
\[
 \left|\sum_v m_v^2-\sum_v\widetilde m_v^2\right|
 \le2\sum_v\widetilde m_v+d=2N+O(R^2)=o(\eta Q).
\]
Here $R^2/N=O(N/Q)=o(1)$, so $d=O(N)$. Thus $\sum m_v^2\le Q$, its deficit is $O(\eta Q)$, and $Z'=(\sum m_v^2-N)/2$ has the stated properties.

For the energy before rounding,
\begin{align*}
 \sum_v\widetilde m_v|v|^2
 &=c(R^2P_1-P_2)
 =N\left(R^2-\frac{P_2}{P_1}\right)\\
 &=\frac13NR^2(1+o(1))
 =(1+o(1))\frac{4N^3}{27Q\Psi_*}.
\end{align*}
The rounding changes it by at most $dR^2=O(R^4)=o(NR^2)$, since $R^2/N=O(N/Q)\to0$. At least two independent primitive directions outside the fixed finite set have multiplicity tending to infinity: on any two fixed allowed directions of row one the normalized radial deficits are bounded away from zero, and the multiplicity scale $cR^2\asymp Q/N$ tends to infinity. This proves rank two. Taking integer coordinates in the displayed lattice basis gives an integral matrix with the same Gram determinant, as the lattice basis has determinant one.
\end{proof}

The trace upper bound implies a Gram-determinant upper bound because $\det M\le(\tr M/2)^2$ for every positive semidefinite $2\times2$ matrix $M$. A gcd greater than one would only decrease the normalized determinant. Thus Lemma~\ref{lem:construction} already proves the sharp upper bound for $E_{\le}$.

\section{Exact integer correction by four squares}
We now prove the upper bound with exactly $z$ zero minors, and in particular show that the exact class is eventually nonempty.
Put
\[
 e=\frac{z}{n^2}+\frac{n}{z},\qquad \ep=e^{1/4},\qquad
 M=\left\lceil\sqrt{\ep z}\right\rceil,
\]
and reserve the eight distinct primitive directions $(u,1)$, $u=0,1,\ldots,7$. They form the excluded set $F$ in Lemma~\ref{lem:construction}. Set
\[
 N=n-8M,\qquad Z=z-4M(M-1),\qquad Q'=N+2Z.
\]
Then $M=o(n)$, $M^2\sim\ep z=o(z)$, $N\sim n$, and $Q'\sim2z\sim Q$. In particular the parameters satisfy $N\ll Q'\ll N^2$ and $Z\ge0$ for large $n$. Apply Lemma~\ref{lem:construction} to $N,Q'$ with $\eta=e^{1/2}$. Its hypothesis holds because $Q'/N^2+N/Q'=O(e)$ and $e/\eta\to0$.
Let $Z'$ be the resulting zero-minor count of the bulk columns. The unused count
\[
 D=Z-Z'\ge0
\]
is an integer with $D=O(e^{1/2}z)=o(M^2)$. Therefore $D\le M^2$ for sufficiently large $n$. By Lagrange's four-square theorem choose integers $x_1,\ldots,x_4$ with
\[
 D=x_1^2+x_2^2+x_3^2+x_4^2,
 \qquad |x_j|\le\sqrt D\le M.
\]
On the reserved pair of directions numbered $2j-2,2j-1$, put multiplicities $M+x_j$ and $M-x_j$. They are nonnegative integers, and all reserved directions are distinct and absent from the bulk. These directions use exactly $8M$ columns and create exactly
\[
 \sum_{j=1}^4\left[\binom{M+x_j}{2}+\binom{M-x_j}{2}\right]
 =4M(M-1)+\sum_{j=1}^4x_j^2
\]
additional zero minors. There are no zero minors between different primitive directions. The combined matrix consequently has $N+8M=n$ columns and
\[
 Z'+4M(M-1)+D=z
\]
zero minors, exactly.

The bulk lattice is rectangular with $a=T_0/R$ and $R^2\asymp N^2/Q'\asymp n^2/z$. Each reserved vector has squared length
\[
 a^2u^2+a^{-2}=O_{T_0}(R^2)
\]
for fixed $0\le u\le7$. The reserved energy is thus $O(MR^2)=o(n^3/Q)$. The bulk energy is $(1+o(1))4N^3/(27Q'\Psi_*)=(1+o(1))4n^3/(27Q\Psi_*)$. Applying the trace upper bound and then dividing by the squared gcd proves
\[
 E_=(n,z)\le(1+o(1))\frac{4n^6}{729Q^2\Psi_*^2}.
\]
This completes Theorem~\ref{thm:main}.

\section{What the theorem says about the construction}
In the bulk construction, write a primitive direction on row $j$ as $(a(u+j\theta),j/a)$ and set $x=T_0(u+j\theta)/R^2$. Up to relative asymptotic normalization and integer rounding, its multiplicity is
\[
 m_{u,j}\simeq\frac{3Q}{2n}\left(1-x^2-\frac{j^2}{T_0^2}\right)_+.
\]
Thus multiplicity decreases quadratically toward the boundary of an ellipse. Only the finitely many rows $j<T_0$ are used in the bulk. The number of distinct directions is of order $R^2\asymp n^2/z$, while the typical multiplicity is of order $z/n$. The exact-count correction uses $o(n)$ columns and an $o(1)$ fraction of the squared-length energy.

Numerical evaluation suggests a global maximizer near $T=10.67368$ and $\Psi_*\approx0.3183946273$, which would give the coefficient $1/(729\Psi_*^2)\approx0.01353135$. These numerical suggestions are not used in the theorem. In particular this paper does not assert uniqueness of the maximizer or a rigorously certified ten-row classification of every optimizer.

\section{Exact certificates and finite illustrations}
The standard-library Python script \path{verify_gram.py} checks the rational certificate of Lemma~\ref{lem:maximum} by integer square roots and rational arithmetic. It gives
\[
 0.318394611593156<\Psi(32/3)<0.318394611593166.
\]
It also constructs explicit primitive integer vectors at $T=32/3$ with multiplicities obtained by rounding the parabolic profile. For the examples below the amplitude and squared radius are both the first column. The determinant and zero-minor count are exact integers. The last column is a decimal rendering of an exact rational ratio; none of these finite matrices is claimed to minimize the determinant.
\begin{center}\small
\begin{tabular}{rrrrr}\toprule
$R^2$ & $n$ & $z$ & directions & $G(A)z^2/n^6$\\\midrule
100 & 4\,841 & 158\,600 & 98 & 0.012378990623\\
300 & 43\,154 & 4\,298\,595 & 284 & 0.013119376466\\
1\,000 & 478\,132 & 159\,219\,258 & 946 & 0.013406349757\\
3\,000 & 4\,299\,968 & 4\,298\,439\,159 & 2\,836 & 0.013489548326\\
10\,000 & 47\,765\,579 & 159\,199\,178\,481 & 9\,466 & 0.013518794709\\\bottomrule
\end{tabular}
\end{center}
An exact correction example uses $1357=0^2+3^2+18^2+32^2$ and $M=40$. The resulting matrix has exactly $18\,048$ columns and $1\,140\,772$ zero minors. The script checks the Gram determinant against the weighted Cauchy--Binet sum, verifies all zero counts, and records the full weighted matrix. Each displayed matrix contains the primitive horizontal direction and a row-one vector, proving that the gcd of its maximal minors is one. Complete output appears in \path{data/verification.json} and \path{data/verification.txt}.

\section{Source relationship and limits of the novelty claim}
The source is ProveIt, \path{Analysis/GeneralizedRiemannHypothesis/Article/sections/src.tex}, pinned at commit \path{cfab41369737c66ae520d179419509eb719367cd}. Its results \texttt{src:thm:gramasym} and \texttt{src:lem:gramcount} develop the primitive-direction potential for the generic rank-two problem, and \texttt{src:q:gram} asks about vanishing minors. The incoming archive \path{gram_column_number_2026-11-27.zip} proves the order $\Theta_r(n^{3r-2})$ for generic fixed-rank matrices and a bound with at most a fixed number of zero minors. It expressly does not determine the exact-zero optimum or its sharp coefficient.

The present theorem addresses a growing-degeneracy regime not covered by either statement. The new ingredients are the $\ell^2$ multiplicity budget, the squared primitive potential, the parabolic construction, and the exact four-square correction. The classical theorem of Lagrange is used only to satisfy the exact integer number of zero minors. The nearby published work of B\'ar\'any--Tokushige minimizes a sum of absolute determinants for convex lattice polygons, while Kriepke--Schymura studies maximum absolute minors and generic column counts. Neither objective is the multiplicity-constrained sum of squared normalized minors treated here. A limited literature search did not locate this theorem; an exhaustive priority claim is not made. No prominent classical conjecture is claimed resolved.

\begin{thebibliography}{9}
\bibitem{proveit} V. Reshetnikov, \emph{ProveIt}, source snapshot above. \url{https://github.com/VladimirReshetnikov/ProveIt/tree/cfab41369737c66ae520d179419509eb719367cd/Analysis/GeneralizedRiemannHypothesis}.
\bibitem{bt} I. B\'ar\'any and N. Tokushige, \emph{The minimum area of convex lattice $n$-gons}, Combinatorica 24 (2004), 171--185. \url{https://www.renyi.hu/~barany/cikkek/94.pdf}.
\bibitem{ks} B. Kriepke and M. Schymura, \emph{On generic $\Delta$-modular integer matrices with two rows}, arXiv:2502.15394v2. Published in SIAM J. Discrete Math. 40 (2026), 1342--1368. \url{https://arxiv.org/abs/2502.15394}; \url{https://doi.org/10.1137/25M1745878}.
\end{thebibliography}
\end{document}
